% XeLaTeX starter: replace metadata and the body between BODY START / BODY END.
\documentclass[UTF8,12pt,a4paper,fontset=none]{ctexart}
\usepackage[a4paper,top=2.5cm,bottom=2.5cm,left=2.7cm,right=2.7cm]{geometry}
\usepackage{graphicx}
\usepackage{longtable,booktabs,array,calc,tabularx}
\usepackage{enumitem}
\usepackage{amsmath,amssymb}
\usepackage{xcolor}
\usepackage{etoolbox}
\usepackage{hyperref}
\usepackage{fancyhdr,titlesec,tocloft,caption,url,indentfirst}
\usepackage[protrusion=false]{microtype}

% CONFIG START: all strings may be replaced for the current report.
\newcommand{\ReportType}{}
\newcommand{\ReportTitle}{中文课程报告版式示例}
\newcommand{\ReportSubtitle}{}
\newcommand{\ReportCollege}{}
\newcommand{\ReportClass}{}
\newcommand{\ReportStudentId}{}
\newcommand{\ReportAuthor}{}
\newcommand{\ReportAdvisor}{}
\newcommand{\ReportHeaderTitle}{中文课程报告版式示例}
\newif\ifShowFigureList
\ShowFigureListtrue
% CONFIG END

% Use regular static faces on Windows; VF files can resolve to ExtraLight/Thin.
\IfFileExists{C:/Windows/Fonts/simsun.ttc}{
  \setmainfont{SimSun}[BoldFont=SimHei]
  \setCJKmainfont{SimSun}[BoldFont=SimHei]
}{\IfFontExistsTF{Noto Serif CJK SC}{
  \setmainfont{Noto Serif CJK SC}\setCJKmainfont{Noto Serif CJK SC}
}{\IfFontExistsTF{Noto Serif SC}{
  \setmainfont{Noto Serif SC}[AutoFakeBold=2]\setCJKmainfont{Noto Serif SC}[AutoFakeBold=2]
}{\IfFontExistsTF{SimSun}{
  \setmainfont{SimSun}[BoldFont=SimHei]\setCJKmainfont{SimSun}[BoldFont=SimHei]
}{
  \setmainfont{FandolSong-Regular}[Extension=.otf,BoldFont=FandolSong-Bold]
  \setCJKmainfont{FandolSong-Regular}[Extension=.otf,BoldFont=FandolSong-Bold]
}}}}
\IfFileExists{C:/Windows/Fonts/Noto Sans SC Bold (TrueType).otf}{
  \setsansfont{Noto Sans SC (TrueType).otf}[Path=C:/Windows/Fonts/,BoldFont={Noto Sans SC Bold (TrueType).otf}]
  \setCJKsansfont{Noto Sans SC (TrueType).otf}[Path=C:/Windows/Fonts/,BoldFont={Noto Sans SC Bold (TrueType).otf}]
}{\IfFontExistsTF{Noto Sans CJK SC}{
  \setsansfont{Noto Sans CJK SC}\setCJKsansfont{Noto Sans CJK SC}
}{\IfFontExistsTF{Noto Sans SC}{
  \setsansfont{Noto Sans SC}\setCJKsansfont{Noto Sans SC}
}{\IfFontExistsTF{SimHei}{
  \setsansfont{SimHei}\setCJKsansfont{SimHei}
}{
  \setsansfont{FandolHei-Regular}[Extension=.otf,BoldFont=FandolHei-Bold]
  \setCJKsansfont{FandolHei-Regular}[Extension=.otf,BoldFont=FandolHei-Bold]
}}}}

% Compatibility verified for LaTeX 2025-11 + longtable 2025-10-13 v4.24.
% Limit this repair to that package date; later releases keep their own output.
\makeatletter
\IfFormatAtLeastTF{2025-11-01}{
  \@ifpackagelater{longtable}{2025/10/13}{
    \@ifpackagelater{longtable}{2025/10/14}{}{
      \patchcmd{\LT@output}{\vss}{\vfil}{}{\PackageWarning{csureport}{Inspect longtable compatibility}}
      \patchcmd{\LT@output}{\vss}{\vfil}{}{\PackageWarning{csureport}{Inspect longtable compatibility}}
    }
  }{}
}{}
\makeatother

\hypersetup{colorlinks=false,hidelinks,pdfauthor={\ReportAuthor},pdftitle={\ReportTitle},bookmarksnumbered=true}
\urlstyle{same}
\setlength{\parindent}{2em}
\setlength{\parskip}{0pt}
% Preserve ctexart's reference spacing: nominal 20pt becomes 26pt at 1.3.
\linespread{1.3}
\setlist{nosep,leftmargin=2em}
\titleformat{\section}{\bfseries\sffamily\fontsize{15pt}{20pt}\selectfont}{\thesection.}{0.6em}{}
\titleformat{\subsection}{\bfseries\sffamily\fontsize{14pt}{20pt}\selectfont}{\thesubsection}{0.6em}{}
\titleformat{\subsubsection}{\bfseries\sffamily\fontsize{12pt}{20pt}\selectfont}{\thesubsubsection}{0.6em}{}
\titlespacing*{\section}{0pt}{18pt}{8pt}
\titlespacing*{\subsection}{0pt}{14pt}{6pt}
\titlespacing*{\subsubsection}{0pt}{10pt}{4pt}
% Reserve space before hyperref places heading anchors at a page boundary.
\newcommand{\ReportHeadingSpace}[1]{\par
  \ifdim\dimexpr\pagegoal-\pagetotal\relax<#1\baselineskip\newpage\fi}
\pretocmd{\section}{\ReportHeadingSpace{4}}{}{\PackageWarning{csureport}{Inspect section spacing}}
\pretocmd{\subsection}{\ReportHeadingSpace{3}}{}{\PackageWarning{csureport}{Inspect subsection spacing}}
\pretocmd{\subsubsection}{\ReportHeadingSpace{3}}{}{\PackageWarning{csureport}{Inspect subsubsection spacing}}
\pagestyle{fancy}
\fancyhf{}
\setlength{\headheight}{39pt}
\newsavebox{\ReportHeaderBox}
\newcommand{\ReportHeaderLine}[1]{\begingroup
  \sbox{\ReportHeaderBox}{#1}%
  \ifdim\wd\ReportHeaderBox>\linewidth
    \resizebox{\linewidth}{!}{\usebox{\ReportHeaderBox}}%
  \else\usebox{\ReportHeaderBox}\fi\endgroup}
\fancyhead[L]{\includegraphics[height=0.95cm]{media/csu_seal.png}}
\fancyhead[C]{\parbox[c][1.0cm][c]{0.48\textwidth}{\centering\sffamily\small 中南大学报告\\[-0.15em]\footnotesize \ReportHeaderTitle}}
\fancyhead[R]{\parbox[c][1.0cm][c]{0.24\textwidth}{\raggedleft\sffamily\footnotesize \ReportHeaderLine{\ReportCollege}\\\ReportHeaderLine{\ReportAuthor\quad\ReportStudentId}}}
\fancyfoot[C]{第\thepage 页}
\renewcommand{\headrulewidth}{0.6pt}
\newcommand{\infoblank}[1]{\underline{\makebox[6.4cm][c]{#1}}}
\newcommand{\ReportCode}[1]{\mbox{#1}}
\captionsetup{font=small,labelfont=bf}
\renewcommand{\contentsname}{目\quad 录}
\renewcommand{\cftsecleader}{\cftdotfill{\cftdotsep}}
\renewcommand{\listfigurename}{图\quad 录}
\tocloftpagestyle{fancy}
\renewcommand{\figurename}{图}
\renewcommand{\tablename}{表}
\setcounter{tocdepth}{2}
\setcounter{secnumdepth}{3}
\sloppy

\begin{document}
\fontsize{12pt}{20pt}\selectfont
\hypersetup{pageanchor=false}
\begin{titlepage}
\thispagestyle{empty}
\begin{center}
\vspace*{0.3cm}
\includegraphics[width=0.92\textwidth]{media/csu_logo_full.png}\\[1.2cm]
\if\relax\detokenize\expandafter{\ReportType}\relax\else
{\sffamily\bfseries\fontsize{20pt}{30pt}\selectfont \ReportType\par}
\vspace{1.6cm}
\fi
{\sffamily\bfseries\fontsize{19pt}{31pt}\selectfont \ReportTitle\par}
\if\relax\detokenize\expandafter{\ReportSubtitle}\relax\else
\vspace{0.9cm}
{\fontsize{14pt}{24pt}\selectfont \ReportSubtitle\par}
\fi
\vfill
\renewcommand{\arraystretch}{1.75}
\begin{tabular}{>{\bfseries\sffamily}r@{\hspace{0.55cm}}l}
学院： & \infoblank{\ReportCollege}\\
班级： & \infoblank{\ReportClass}\\
学号： & \infoblank{\ReportStudentId}\\
姓名： & \infoblank{\ReportAuthor}\\
指导教师： & \infoblank{\ReportAdvisor}\\
\end{tabular}
\vfill
{\sffamily\fontsize{12pt}{20pt}\selectfont CENTRAL SOUTH UNIVERSITY\par}
\end{center}
\end{titlepage}
\hypersetup{pageanchor=true}
\pagenumbering{Roman}
\tableofcontents
\clearpage
\ifShowFigureList
\listoffigures
\clearpage
\fi
\pagenumbering{arabic}

% BODY START: replace this illustrative content with the current assignment.
\begin{center}
{\sffamily\bfseries\fontsize{18pt}{28pt}\selectfont \ReportTitle\par}
\end{center}
\vspace{0.9em}
\begin{center}
{\sffamily\bfseries\fontsize{14pt}{20pt}\selectfont 摘\quad 要}
\end{center}
本文件展示封面、正文标题、目录、图录及图表的排版。正式报告的摘要应说明研究问题、材料与方法、主要结果和适用范围，以下示例内容应随实际任务整体替换。

\textbf{关键词：}课程报告；中文排版；图表规范

\section{报告背景}
正文采用十二磅中文衬线字，首行缩进两字。底稿保留参考源码的二十磅名义行距与ctex的1.3倍行距设置，实际基线间距为二十六磅。章节名称和内容按实际报告组织。
\subsection{资料与方法}
资料来源、研究范围与处理方法应与报告内容对应。
\subsubsection{资料整理}
使用三级编号组织必要的小节，目录仅显示前两级标题。

\section{图表与结果}
\subsection{表格示例}
\begin{longtable}{@{}p{0.16\textwidth}p{0.32\textwidth}p{\dimexpr0.52\textwidth-4\tabcolsep\relax}@{}}
\toprule
\textbf{编号} & \textbf{项目} & \textbf{说明}\\
\midrule
\endfirsthead
\toprule
\textbf{编号} & \textbf{项目} & \textbf{说明}\\
\midrule
\endhead
\bottomrule
\endlastfoot
\ReportCode{UC-01} & 短编号 & 编号保持整行，长说明正常换行。\\
\ReportCode{BR-01} & 跨页表 & 表头在每一续页重复，表格保留三线结构。\\
\end{longtable}

\subsection{图形示例}
图\ref{fig:layout-demo}展示图号与题注位置。正式图应按竖向页面重新布局，保留可读文字与完整标记。
\begin{figure}[htbp]
\centering
\fbox{\parbox[c][4cm][c]{0.84\textwidth}{\centering 黑白示意图位置\\[0.4em]替换为当前报告的实际图形}}
\caption{图形与题注排版示例}
\label{fig:layout-demo}
\end{figure}

\section*{参考文献}
\phantomsection
\addcontentsline{toc}{section}{参考文献}
\noindent 示例条目仅展示排版。正式内容应写明作者或机构、标题、发布日期、网址及访问日期，并让每个条目保持完整。

\appendix
\section{补充材料}
附录沿用正文的字体、页面和图表规范，内容依据实际任务提供。
% BODY END
\end{document}
